% =====================================================================================================
% Online Appendix F: The Revealed Value of Compactness: Robustness (v2).
% Writer appendix_D_F, 2026-09-24. Sources: output/memos/ra2_wedge.md (+ review), ra1_modelfree.md (Farseer),
% ra5_theory.md (wedge cases, partial identification), ra3_econ.md (aggregates); output/tables/ra2_wedge_*.csv,
% ra1_modelfree_farseer_*.csv, ra5_theory_pi_*.csv, ra3_econ_inference_share*.csv. Tables: paper/tables/appF_*.tex.
% Numbers and sources: paper/notes/writer2_appendix_D_F.md.
% =====================================================================================================

\section{The Revealed Value of Compactness: Robustness}\label{app:wedge}
\setlength{\emergencystretch}{2em}% allow looser lines in the online appendix

This appendix supports Section~\ref{sec:wedge}. Appendix~\ref{app:wedge-cases} states what the wedge identifies under
alternative objectives of the developer. Appendices~\ref{app:wedge-sample} through~\ref{app:wedge-labown} document the
clean inference-demand sample, the 32 technologies, parameter-count conventions and the developers' own technologies.
Appendix~\ref{app:wedge-pi} treats extrapolation beyond the designs, and Appendices~\ref{app:wedge-conduct}
through~\ref{app:wedge-trends} report the conduct tests, family token budgets, the cost side, validation and trends.
Unless stated otherwise, wedges are evaluated under the reference technology, the Chinchilla form with the outer
exponent free ($\sigma^*_\kappa=0.701$), on the 77 models of the clean inference-demand sample, and intervals for sample
statistics come from the joint distribution across bootstrap draws of the technology: each draw recomputes the median
and the shares for the whole sample.

% -----------------------------------------------------------------------------------------------------
\subsection{What the Wedge Identifies}\label{app:wedge-cases}

Table~\ref{tab:app-wedge-cases} lists what the wedge $w=\varepsilon_N/\varepsilon_D$ and the share $s=(w-1)/w$ measure
under the objectives and cost structures of Proposition~\ref{prop:A-wedge}. The rows fall into three groups. In the
first, a developer bears some or all of the cost of serving its model, and $w-1=\lambda\phi X_S/X_T$ is internalized
serving expenditure relative to training expenditure. This ratio needs no conversion of FLOPs into dollars; a count of
serving tokens needs the price ratio $p$ as well. In the second group, a binding memory or tier cap or a latency
objective multiplies the wedge by one plus the constraint's normalized multiplier. The wedge then bounds the value of
compactness from above without identifying it, and any wedge can be rationalized with no serving demand at all. In the
third group, data costs, data caps and logit distillation add to the cost of tokens, so the wedge understates the value
of compactness; a binding data cap can push it below one. Serving demand alone cannot produce $w<1$ as long as serving
cost rises with size at given quality.

The last row is the factor-bias critique of markup estimation. If the first-order conditions of the developer involve a
technology or an output different from the one the econometrician uses, the estimated wedge is the true one times the
ratio of the two technologies' wedges at the same inputs. The developer's own technology (Appendix~\ref{app:wedge-labown})
and the spread across 32 technologies (Appendix~\ref{app:wedge-tech}) are the two ways the paper addresses it. The cases
also make different predictions about conduct (Remark~\ref{rem:A-conduct}), which Appendix~\ref{app:wedge-conduct} tests.

% Online Appendix table: what the wedge identifies (output/tables/ra5_theory_wedge_cases.tex), relabelled. Writer appendix_D_F.
% v3 (writer appendices): prices q_T, q_D; latency Z(N); binding data cap for a member (d') and a family (g'), post-training;
% the family and cap rows verified by code/paper/verify_family_cap.py.
% Round 3 (WP3, fix list P1(j), P10(i), P10(k)): the value of compactness is the virtual value v_N = m_N + nu_N; size caps,
% latency objectives and size menus identify v_N and bound m_N from above; token caps bound v_N from below; distillation
% with an unchanged student technology; post-training with m_N net of X_P; notation varrho, h_i, mathcal A. Virtual-value
% rows verified by code/analysis/rb5_theory/verify_virtual_value.py.
% [integration, 2026-09-25] Split into two floats (rows to (c); rows (d) on with the notes, headed "(continued)"),
% because the single float was 80pt taller than the page. Content unchanged.
\begin{table}[tp]
\centering
\caption{What the Wedge Identifies under Alternative Developer Objectives}
\label{tab:app-wedge-cases}
\footnotesize
\setlength{\tabcolsep}{3pt}
\begin{tabular}{@{}>{\raggedright\arraybackslash}p{0.25\textwidth} >{\raggedright\arraybackslash}p{0.275\textwidth} >{\raggedright\arraybackslash}p{0.265\textwidth} >{\raggedright\arraybackslash}p{0.13\textwidth}@{}}
\toprule
Developer objective or cost structure & Wedge $w=\varepsilon_N/\varepsilon_D$ & $s=(w-1)/w$ measures & $\hat T$ vs.\ $\lambda\varrho T$ \\
\midrule
Lifetime compute, developer serves ($\lambda=\varrho=\phi=1$) & $1+T/(3D)$ & Serving share of lifetime compute & Equal \\
(a) Serving cost borne at rate $\lambda$, price ratio $\varrho$, cost elasticity $\phi$ & $1+\lambda\phi\varrho\,T/(3D)$ & $\lambda\phi X_S/(X_T+\lambda\phi X_S)$ & $\times\phi$ \\
(a$'$) Open release, adoption value $v(\mathcal A)$, users bear serving & $1+v'\mathcal A|\varepsilon_{\mathcal A,P}|/X_T$ & Adoption value of compactness & Not $T$ \\
(b) Binding memory or tier cap, shadow value $\nu_N\ge0$ & $1+v_N=1+m_N+\nu_N$ & Identifies $v_N/(1+v_N)$; upper bound on $m_N/(1+m_N)$ & Overstates \\
(b$'$) Menu of sizes, shadow value $\nu_N$ at the chosen size & $1+v_N=1+m_N+\nu_N$ & Identifies $v_N/(1+v_N)$; upper bound on $m_N/(1+m_N)$ if the menu binds as a cap & Overstates if it binds as a cap \\
(c) Binding latency objective, $Z(N)\le\bar Z$ & $1+v_N$, $\nu_N=\ell\xi_Z/X_T$ & Identifies $v_N/(1+v_N)$; upper bound on $m_N/(1+m_N)$ & Overstates \\
\bottomrule
\end{tabular}
\end{table}

\begin{table}[tp]
\centering
{\captionstyle\captionsize Table~\ref{tab:app-wedge-cases} (continued). What the Wedge Identifies under Alternative Developer
Objectives\par}
\vspace{7pt}
\footnotesize
\setlength{\tabcolsep}{3pt}
\begin{tabular}{@{}>{\raggedright\arraybackslash}p{0.25\textwidth} >{\raggedright\arraybackslash}p{0.275\textwidth} >{\raggedright\arraybackslash}p{0.265\textwidth} >{\raggedright\arraybackslash}p{0.13\textwidth}@{}}
\toprule
Developer objective or cost structure & Wedge $w=\varepsilon_N/\varepsilon_D$ & $s=(w-1)/w$ measures & $\hat T$ vs.\ $\lambda\varrho T$ \\
\midrule
(d) Data cost $q_D$ per token & $(1+v_N)/(1+m_D)$ & $(X_S-X_D)/(X_T+X_S)$ & Understates \\
(d$'$) Binding data cap, multiplier $\mu$ & $(1+v_N)/(1+m_D+\mu)$ & Lower bound on $v_N/(1+v_N)$; $w<1$ possible & Understates \\
(e) Logit distillation, teacher $N_T$, student technology unchanged & $\dfrac{1+T/(3D)}{1+N_T/(3N)}$ & $(X_S-X_{\rm teach})/(X_T+X_S)$ & Understates \\
(f) Dollar cost $\propto N^{1+\delta}D$; serving $\propto N^{1+\delta_S}$ & $1+\delta+\lambda(1+\delta_S)\varrho T/(3D)$ & Mixes $\delta$ with serving & $+3D\delta+\lambda\delta_S\varrho T$ \\
(g) Family token budget, sizes chosen or set by binding size constraints & $\sum_i\pi_i(1+v_{N,i})=(1+m_D)\bar w_H$ & $1-1/\bar w_H=\sum_ih_is_i$ & Only $\sum_i\pi_i\hat T_i$ identified \\
(g$'$) Family budget at a binding corpus cap & $\sum_i\pi_i(1+v_{N,i})=(1+m_D+\mu)\bar w_H$ & Family share a lower bound; with $m_D=0$, member $w_i-1\le v_{N,i}$ if and only if the member would train on more tokens at its own size & Understates \\
Post-training compute $x_P=X_P/X_T$ & $1+x_P+m_N$ ($m_N$ net of $X_P$) & Mixes $x_P$ with $m_N$ & Overstates \\
Lab's technology differs (factor bias, beliefs, output) & $\hat w=w_{L'}\,w_L/w_{L'}$ & Contaminated & Over iff $w_L>w_{L'}$ \\
\bottomrule
\end{tabular}
\par\smallskip
\begin{tablenotes}The developer chooses $(N,D)$ to maximize $V(L,N)$ minus training expenditure $X_T=q_T\,6ND$ (and data cost $q_DD$). $m_N\equiv-(\partial V/\partial\ln N)_L/X_T$ is the objective's own value of compactness at fixed quality per unit of training expenditure; $\nu_N$ is the shadow value of binding constraints on size alone (memory or tier caps, latency objectives, menus of sizes) per unit of training expenditure, nonnegative for a cap and negative for a menu that binds as a floor; $v_N\equiv m_N+\nu_N$ is the value of compactness, which equals $m_N$ when no size constraint binds; rows other than (d), (d$'$), (e) and (g)--(g$'$) take $m_D=0$, and rows other than (f) take $\delta=0$. $m_D\equiv q_DD/X_T$; $X_S$ is serving expenditure, $\lambda$ the share of it the developer internalizes, $\varrho$ the price of a serving FLOP relative to a training FLOP, $\phi$ the size elasticity of serving cost; $\varepsilon_{\mathcal A,P}$ is the elasticity of adoption with respect to users' serving cost; $\xi_Z$ and $\mu\ge0$ are the multipliers of the latency objective and the data cap, $\ell$ the size elasticity of latency; in the post-training row, post-training expenditure $X_P$ is counted as a cost alongside $X_T$, so that $m_N$ is net of it; $\delta$ ($\delta_S$) is the size elasticity of the price of a training (serving) FLOP; in rows (g) and (g$'$), $v_{N,i}$ is member $i$'s value of compactness, $\pi_i\propto h_i/w_i$ with $h_i$ the member's compute share, and $\bar w_H$ the compute-weighted harmonic mean of member wedges; in the last row, $w_{L'}$ is the wedge that the lab's technology or output assigns to its choice and $w_L$ the one that the econometrician's assigns to the same inputs. The last column compares $\hat T\equiv3D(\hat w-1)$ with $\lambda\varrho T$, the internalized serving demand in training-FLOP units (rows other than (a) and (f) take $\phi=1$); ``Over iff $w_L>w_{L'}$'' means that $\hat T$ overstates, which under factor bias is $\chi<0$. If distillation also raises the productivity of tokens ($\chi>0$), it lowers the measured wedge further. Every row was verified by brute-force maximization of the objective (maximum relative error $2\times10^{-7}$; Propositions~\ref{prop:A-wedge} and~\ref{prop:A-family}).\end{tablenotes}
\begin{tablenotes}[Source]Author's derivations.\end{tablenotes}
\end{table}


% -----------------------------------------------------------------------------------------------------
\subsection{The Clean Inference-Demand Sample}\label{app:wedge-sample}

The clean sample was defined before any wedge was computed. It starts from the 173 general-purpose open-weight base
models of the verified sample, whose parameter and token counts are documented by their developers (Online
Appendix~\ref{app:data}). Nine steps, applied in order, drop models whose $(N,D)$ is not the choice that
Proposition~\ref{prop:wedge} describes; a model is dropped at the first step that applies. The replication package lists
every flag with its reason and source.
\begin{itemize}
\item[(a)] \emph{One observation per pretraining run} (8 models; the first release is kept): Llama~3.1 8B and 70B, which
continue the pretraining of the Llama~3 runs; Qwen1.5-72B, with the same $N$ and $D$ as Qwen-72B; OpenLLaMA v2 3B and
7B; Yi-1.5 6B and 34B, which the model card describes as continued pretraining of Yi on 500 billion tokens; and
H2O-Danube2-1.8B, initialized from H2O-Danube-1.8B.
\item[(b1)] \emph{Research suites whose $D$ is fixed by design} (45): Pythia, OPT, BLOOM, RWKV-4-Pile, XGLM,
Cerebras-GPT and GPT-Neo, GPT-J and GPT-NeoX; the natural-language stage of CodeGen; and Falcon-RW-1B, a data ablation.
\item[(b2)] \emph{Replications and data benchmarks} (8): OpenLLaMA, RedPajama-INCITE, LLM360 Amber and K2, and DCLM-7B.
\item[(c)] \emph{Mixture-of-experts models} (12), reported with bounds in Appendix~\ref{app:wedge-conv}.
\item[(d)] \emph{Distilled, pruned or derived initialization} (10): Gemma~2 2B and 9B and every Gemma~3 size, which
their technical reports describe as trained with knowledge distillation; Llama~3.2 1B and 3B, pruned from Llama~3.1 8B
and distilled from larger models' logits; and Yi-1.5-9B, continued from a depth-upscaled Yi-6B. A keyword search of all
179 saved model cards and of these reports found no other case.
\item[(e)] \emph{Multimodal pretraining}: Gemma~3 4B--27B and Llama~4, all already removed at (c) or (d).
\item[(f)] \emph{No base checkpoint} (7): Phi-3 mini, small and medium, Phi-3.5-mini, Phi-4-mini, phi-4, and Qwen3-32B.
\item[(g)] \emph{Non-transformer or hybrid architectures}: RWKV, already removed at (b1).
\item[(h)] \emph{Token counts without a primary source} (6): Qwen1.5 0.5B to 32B, whose counts appear only in the
ObsScaling file \citep{ruan2024observational}.
\end{itemize}

The result is 77 models from 36 families and 18 developers, released 2023--2025. Table~\ref{tab:app-cleaning}, panel~A,
follows the headline through the steps. The median wedge stays between 3.77 and 4.46: dropping the research suites and
replications, whose wedges are lower, raises it, and dropping mixture-of-experts, distilled and instruction-only models
lowers it again. The sign is less robust than the point estimates suggest: every one
of the 32 technologies puts the point estimate above one for 64 percent of clean models, but the union of their 95 percent
intervals lies above one for only 18 percent.

Panel~B splits the clean sample. Families that train every size on one token budget, whose member wedges are not revealed
preferences (Appendix~\ref{app:wedge-family}), have a median wedge of 3.90, against 4.01 for families with size-specific
budgets; without them the median share is 0.75. Dropping the nine models trained mainly on teacher-generated synthetic
data (phi-1.5, phi-2 and SmolLM 1--3) gives a median share of 0.74. Models of developers that served their own language
models at release have a median share of 0.80 and a median wedge of 4.95, against 0.74 and 3.80 for the others; the 12
models with a stated on-device or local target, from three developers, have 0.86. The last rows compare models inside and
outside the ranges of $M$ that designs observe. Inside Chinchilla's range ($M\le341$) the median share is 0.56 and no model
has all technologies' intervals above one; only 2 of the 77 models lie inside Chinchilla's design in both $M$ and compute.

% Online Appendix Table E2 (round 3, WP4b, fix list W24(c)): copy of output/tables/rb5_units_appE_cleaning_body.tex,
% generated by code/analysis/rb5_units/paper_tables.py. Panel A: output/tables/rb2_decisions_cleaning.csv (audited token
% counts; rb2_decisions re-run, reviewed). Panel B: per-model wedges of the 32 technologies (exhibits_cache.pkl, L) with the
% budget-level unit kinds of data/processed/rb5_units/units_primary.csv; the generator asserts that the clean-sample and
% version-3 synthetic rows reproduce rb2_decisions_cleaning.csv. Numbers: output/tables/rb5_units_appE_cleaning_panelB.csv.
\begin{table}[tp]
\centering
\caption{The Clean Sample: Construction and Subsamples}
\label{tab:app-cleaning}
\footnotesize
\setlength{\tabcolsep}{2pt}
\begin{tabular*}{\textwidth}{@{\extracolsep{\fill}}lrrcccccc@{}}
\toprule
\multicolumn{9}{@{}l}{\textit{Panel A. Sequential cleaning steps (reference technology: Chinchilla, $\kappa$ free)}} \\[1pt]
 & & & Median & Share & All & Band & Median & Median $w$, \\
Step (models dropped) & $n$ & Groups & $w$ & $w>1$ & points $>1$ & $>1$ & $s$ & $\kappa=1$ \\
\midrule
Verified sample & 173 & 67 & 3.50 & 0.93 & 0.61 & 0.17 & 0.71 & 3.09 \\
(a) One per pretraining run ($-8$) & 165 & 65 & 3.77 & 0.93 & 0.61 & 0.18 & 0.73 & 3.19 \\
(b1) Research suites ($-45$) & 120 & 56 & 4.20 & 0.98 & 0.69 & 0.16 & 0.76 & 3.51 \\
(b2) Replications, benchmarks ($-8$) & 112 & 51 & 4.46 & 0.98 & 0.71 & 0.17 & 0.78 & 3.70 \\
(c) Mixture of experts ($-12$) & 100 & 43 & 4.30 & 0.98 & 0.68 & 0.18 & 0.77 & 3.61 \\
(d) Distilled, pruned, derived ($-10$) & 90 & 40 & 4.03 & 0.98 & 0.64 & 0.17 & 0.75 & 3.42 \\
(e) Multimodal ($-0$) & 90 & 40 & 4.03 & 0.98 & 0.64 & 0.17 & 0.75 & 3.42 \\
(f) No base checkpoint ($-7$) & 83 & 37 & 3.99 & 0.98 & 0.63 & 0.18 & 0.75 & 3.34 \\
(g) Non-transformer ($-0$) & 83 & 37 & 3.99 & 0.98 & 0.63 & 0.18 & 0.75 & 3.34 \\
(h) $D$ undocumented ($-6$): clean & 77 & 36 & 3.99 & 0.97 & 0.62 & 0.18 & 0.75 & 3.34 \\
\addlinespace
\multicolumn{9}{@{}l}{\textit{Panel B. Subsamples of the clean sample}} \\[1pt]
Common-budget members & 45 & 16 & 3.79 & 0.98 & 0.53 & 0.18 & 0.74 & \\
Size-specific members & 17 & 9 & 4.42 & 0.94 & 0.76 & 0.29 & 0.77 & \\
Single-model releases & 15 & 15 & 3.93 & 1.00 & 0.73 & 0.07 & 0.75 & \\
Decisions of one model & 32 & 24 & 4.03 & 0.97 & 0.75 & 0.19 & 0.75 & \\
No synthetic data, earlier flag & 68 & 31 & 3.80 & 0.97 & 0.59 & 0.15 & 0.74 & \\
No synthetic data, protocol rule$^b$ & 39 & 22 & 2.66 & 0.95 & 0.41 & 0.05 & 0.62 & \\
This size served by its developer$^a$ & 26 & 6 & 5.84 & 1.00 &  &  & 0.83 & \\
Not served & 51 & 15 & 3.80 & 0.96 &  &  & 0.74 & \\
On-device target & 9 & 2 & 9.54 & 1.00 &  &  & 0.90 & \\
In a memory-tier window & 49 & 15 & 3.11 & 0.98 &  &  & 0.68 & \\
Outside the tier windows & 28 & 10 & 7.08 & 0.96 &  &  & 0.86 & \\
$M\le341$ (Chinchilla's range) & 31 &  & 2.26 & 0.94 &  &  & 0.56 & \\
$M>341$ & 46 &  & 6.64 & 1.00 &  &  & 0.85 & \\
\bottomrule
\end{tabular*}
\par\smallskip
\begin{tablenotes}
Panel A: sequential steps from the 173 verified open-weight base models, with the audited token counts (StableLM-Alpha
at 0.8 trillion tokens and MPT-30B at 1.05 trillion; Online Appendix~\ref{app:choices}); each row describes the sample
that remains after the step, and a model is dropped at the first step that applies. Reference technology: the Chinchilla
form with the outer exponent free (inner exponents 0.424 and 0.431, $\sigma^*_\kappa=0.701$), in total parameters.
$s=(w-1)/w$. All points $>1$: every one of the 32 technologies of Table~\ref{tab:app-technologies} puts the point
estimate above one. Band $>1$: the union of the 95 percent intervals over those technologies lies above one. With the
token counts of the previous version the all-points share of the clean sample was 0.64; the StableLM-Alpha correction puts
StableLM-Alpha 3B below one under MiniCPM's law, which lowers it to 0.62. Last column: the median under the Chinchilla
form with $\kappa=1$. Groups: families, except developers in the serving, target and tier rows of Panel~B. Panel~B, first
four rows: members of a family whose token budgets are within 10 percent of each other form one decision (Online
Appendix~\ref{app:wedge-family}). Synthetic data: the earlier flag drops the nine models of phi-1.5, phi-2 and SmolLM 1--3;
$^b$the protocol's rule flags every decision whose report documents teacher-generated data in pretraining or mid-training
(21 of the 49 decisions, with 38 models, including all six decisions of 2025). $^a$A first-party per-token API or consumer
product offered this size within 180 days of release (Online Appendix~\ref{app:choices}). Tier windows: 2.4--3.3,
6.5--9.5, 11.5--14.9, 26--32.9 and 65--72.9 billion parameters, where a binding tier's shadow value is part of the value
of compactness.
\end{tablenotes}
\begin{tablenotes}[Source]
Model cards, technical reports and Hugging Face configuration files; author's calculations.
\end{tablenotes}
\end{table}


% -----------------------------------------------------------------------------------------------------
\subsection{Technologies}\label{app:wedge-tech}

\emph{The ex-ante rule.}---The set of technologies was fixed before any wedge was computed. It contains (i) every
technology estimated in Section~\ref{sec:tech} and Online Appendix~\ref{app:additional} from final-checkpoint losses of
a designed sweep with the primary estimator (Huber loss on log loss), with each corpus and parameter-count convention as
its own row, under $\kappa=1$ and, where estimated, with $\kappa$ free; (ii) for each IsoFLOP design, the Approach-2 path
with its $\kappa=1$ curvature, the primal fit, and the path combined with the design's model-free curvature; and (iii)
the literature's estimates on the Chinchilla data and the allocation laws published by Meta, DeepSeek
\citep{bi2024deepseek} and MiniCPM \citep{hu2024minicpm}. Excluded ex ante are DataDecide and the checkpoint-interpolated
IsoFLOPs of \citet{li2025misfitting}, whose losses come before the learning-rate schedule ends, an artifact that varies
with $D/N$, and estimator or sample variants of a design already in the set, which are reported below as sensitivity
checks. Table~\ref{tab:app-technologies} lists the 32 technologies.

When a path and a curvature come from different sources, the inner exponents are $\alpha=S(1-a)$ and $\beta=Sa$ with
$S=2(1/\sigma^*-1)$; the wedge depends on the technology only through $S$, the path slope $a$ and the path level
(Corollary~\ref{cor:suff}). DeepSeek's law is stated in non-embedding FLOPs per token, which we compute for other
developers' models from their configuration files; MiniCPM's in non-embedding parameters. MiniCPM's compute-optimal ratio
of 192 tokens per parameter at $10^{21}$ FLOP lies outside its own design, whose ratios run from 10 to 60, so it is itself
an extrapolation. The reference technology and the model-free rows are bootstrapped design-conditionally (wild bootstrap
with the weights of \citealt{feng2011wild}); the other rows reuse the pairs or stratified draws of their source
estimates, and published laws enter as points.

\emph{What drives the dispersion.}---The median share ranges from 0.17 (MiniCPM) to 0.85 (Muennighoff et al.'s runs
with $\kappa$ free). Technologies that let the data choose the curvature give 0.69 to 0.79. The low medians come from
three sources: $\kappa=1$ curvatures far above the model-free estimates, as in Marin's Approach-2 laws ($\sigma^*$ of 0.87
to 0.90, from designs spanning two decades of compute on which the frontier's $E$ and $\gamma$ are not separately
identified); the OLMo ladder with $\kappa=1$, whose $M^*(10^{24})$ of 291 puts a quarter of the sample below one; and
MiniCPM's zero point. The high medians come from the $\kappa$-free fits of the small sweeps, whose curvatures (0.51 to
0.62) are heterogeneous.

Curvature governs levels, not signs. Holding the reference zero point fixed, lowering $\sigma^*$ from 0.80 to 0.60 raises
the median wedge from 2.25 to 8.64 and the median share from 0.55 to 0.88, and Llama~3 8B's wedge from 3.0 to 19.5. The
share of models with $w>1$, 0.974, does not move. At $\sigma^*=0.70$ the median wedge is 4.00, almost exactly the
reference value.

\emph{Estimator variants.}---On the Chinchilla data, least squares in levels gives a median wedge of 4.58, keeping the
five outliers gives 5.50, and the Chinchilla form with $\kappa=1$ gives 3.34 (Table~\ref{tab:app-technologies}). The
bootstrap matters for intervals, not for points. For the reference, pairs resampling, which changes the design, gives
a 95 percent interval for the median wedge of [3.25, 6.04], against [3.21, 4.86] from the design-conditional wild
bootstrap; for the $\kappa=1$ refit, a cluster bootstrap over the nine budgets gives [1.83, 6.48]. A separate
implementation of the $\kappa$-free fit differs from ours in the fourth decimal; over all 173 verified models it gives a
median wedge of 3.78, against 3.77.

% Online Appendix Table F3: all 32 ex-ante technologies on the clean sample (output/tables/ra2_wedge_cleaning.tex Panel B;
% ra2_wedge_technologies.csv). Writer appendix_D_F, 2026-09-24.
\begin{table}[tp]
\centering
\caption{The Revealed Wedge under Every Ex-Ante Technology}
\label{tab:app-technologies}
\footnotesize
\setlength{\tabcolsep}{2.4pt}
\begin{tabular*}{\textwidth}{@{\extracolsep{\fill}}lccr@{\,}lcr@{\,}l@{}}
\toprule
Technology & $\sigma^*$ & $M^*(10^{24})$ & \multicolumn{2}{c}{Median $w$} & Share $w>1$ & \multicolumn{2}{c}{Median $s$} \\
\midrule
\multicolumn{8}{@{}l}{\textit{Chinchilla data}} \\
\quad $\kappa$ free (reference) & 0.701 & 21.6 & 3.99 & {\footnotesize [3.21, 4.86]} & 0.97 & 0.75 & {\footnotesize [0.69, 0.79]} \\
\quad $\kappa=1$ & 0.737 & 17.6 & 3.34 & {\footnotesize [2.64, 4.81]} & 1.00 & 0.70 & {\footnotesize [0.62, 0.79]} \\
\quad Besiroglu et al.\ parameters & 0.737 & 18.1 & 3.31 &  & 1.00 & 0.70 &  \\
\quad Hoffmann et al., Approach 3 & 0.762 & 62.0 & 2.12 &  & 0.88 & 0.53 &  \\
\quad Non-embedding $N$, $\kappa=1$ & 0.753 & 9.7 & 3.65 & {\footnotesize [2.83, 4.71]} & 1.00 & 0.73 & {\footnotesize [0.65, 0.79]} \\
\addlinespace[2pt]
\multicolumn{8}{@{}l}{\textit{Farseer}} \\
\quad Non-embedding $N$, $\kappa=1$ & 0.772 & 17.3 & 2.74 & {\footnotesize [2.38, 3.20]} & 1.00 & 0.63 & {\footnotesize [0.58, 0.69]} \\
\quad Total $N$, $\kappa=1$ & 0.724 & 67.7 & 2.64 & {\footnotesize [2.23, 3.10]} & 0.88 & 0.62 & {\footnotesize [0.55, 0.68]} \\
\quad Non-embedding $N$, $\kappa$ free & 0.710 & 18.3 & 3.91 & {\footnotesize [3.77, 4.06]} & 1.00 & 0.74 & {\footnotesize [0.73, 0.75]} \\
\addlinespace[2pt]
\multicolumn{8}{@{}l}{\textit{Gadre et al.}} \\
\quad RefinedWeb, $\kappa=1$ & 0.796 & 1.4 & 4.41 & {\footnotesize [2.19, 7.67]} & 1.00 & 0.77 & {\footnotesize [0.54, 0.87]} \\
\quad C4, $\kappa=1$ & 0.795 & 6.0 & 3.18 & {\footnotesize [1.12, 7.66]} & 1.00 & 0.69 & {\footnotesize [0.11, 0.87]} \\
\quad RedPajama, $\kappa=1$ & 0.795 & 3.9 & 3.50 & {\footnotesize [2.52, 6.94]} & 1.00 & 0.71 & {\footnotesize [0.60, 0.86]} \\
\quad RefinedWeb, $\kappa$ free & 0.594 & 61.0 & 4.76 & {\footnotesize [3.49, 15.39]} & 0.88 & 0.79 & {\footnotesize [0.71, 0.94]} \\
\quad C4, $\kappa$ free & 0.607 & 78.5 & 3.95 & {\footnotesize [2.06, 10.60]} & 0.88 & 0.75 & {\footnotesize [0.51, 0.91]} \\
\quad RedPajama, $\kappa$ free & 0.623 & 32.3 & 5.64 & {\footnotesize [3.48, 12.88]} & 0.94 & 0.82 & {\footnotesize [0.71, 0.92]} \\
\addlinespace[2pt]
\multicolumn{8}{@{}l}{\textit{OLMo ladder}} \\
\quad $\kappa=1$ & 0.825 & 291 & 1.32 & {\footnotesize [0.45, 2.66]} & 0.77 & 0.24 & {\footnotesize [$-$1.22, 0.62]} \\
\quad $\kappa$ free & 0.544 & 142 & 4.68 & {\footnotesize [3.80, 6.37]} & 0.84 & 0.79 & {\footnotesize [0.74, 0.84]} \\
\addlinespace[2pt]
\multicolumn{8}{@{}l}{\textit{Muennighoff et al.}} \\
\quad $\kappa=1$ & 0.735 & 36.7 & 2.55 & {\footnotesize [0.54, 10.02]} & 0.92 & 0.61 & {\footnotesize [$-$0.85, 0.90]} \\
\quad $\kappa$ free & 0.511 & 100.0 & 6.87 & {\footnotesize [1.15, 37.30]} & 0.86 & 0.85 & {\footnotesize [0.13, 0.97]} \\
\addlinespace[2pt]
\multicolumn{8}{@{}l}{\textit{Meta (Llama 3 IsoFLOPs)}} \\
\quad Meta's law (A2, $\kappa=1$) & 0.785 & 31.3 & 2.30 & {\footnotesize [1.95, 2.65]} & 0.92 & 0.57 & {\footnotesize [0.49, 0.62]} \\
\quad Primal, $\kappa=1$ & 0.769 & 9.6 & 3.26 & {\footnotesize [2.36, 4.04]} & 1.00 & 0.69 & {\footnotesize [0.58, 0.75]} \\
\quad A2 path, model-free $\sigma^*$ & 0.660 & 31.3 & 4.79 & {\footnotesize [3.60, 6.58]} & 0.92 & 0.79 & {\footnotesize [0.72, 0.85]} \\
\addlinespace[2pt]
\multicolumn{8}{@{}l}{\textit{Marin ladders}} \\
\quad Nemotron-CC, primal & 0.658 & 81.8 & 3.52 & {\footnotesize [2.45, 6.84]} & 0.87 & 0.72 & {\footnotesize [0.59, 0.85]} \\
\quad DCLM, primal & 0.647 & 106 & 3.31 & {\footnotesize [2.41, 6.40]} & 0.86 & 0.70 & {\footnotesize [0.59, 0.84]} \\
\quad Comma, primal & 0.663 & 228 & 2.13 & {\footnotesize [1.48, 3.97]} & 0.78 & 0.53 & {\footnotesize [0.32, 0.75]} \\
\quad Nemotron-CC, A2, $\kappa=1$ & 0.867 & 35.1 & 1.63 & {\footnotesize [1.36, 2.17]} & 0.92 & 0.39 & {\footnotesize [0.27, 0.54]} \\
\quad DCLM, A2, $\kappa=1$ & 0.866 & 32.5 & 1.66 & {\footnotesize [1.26, 2.40]} & 0.92 & 0.40 & {\footnotesize [0.21, 0.58]} \\
\quad Comma, A2, $\kappa=1$ & 0.904 & 58.0 & 1.33 & {\footnotesize [1.13, 1.57]} & 0.88 & 0.25 & {\footnotesize [0.11, 0.36]} \\
\quad Nemotron-CC, model-free $\sigma^*$ & 0.705 & 35.1 & 3.78 & {\footnotesize [2.48, 5.79]} & 0.92 & 0.74 & {\footnotesize [0.60, 0.83]} \\
\quad DCLM, model-free $\sigma^*$ & 0.713 & 32.5 & 3.71 & {\footnotesize [2.37, 6.16]} & 0.92 & 0.73 & {\footnotesize [0.58, 0.84]} \\
\quad Comma, model-free $\sigma^*$ & 0.700 & 58.0 & 3.21 & {\footnotesize [2.40, 5.38]} & 0.88 & 0.69 & {\footnotesize [0.58, 0.81]} \\
\addlinespace[2pt]
\multicolumn{8}{@{}l}{\textit{Published laws}} \\
\quad DeepSeek LLM, reference $\sigma^*$ & 0.701 & 13.9 & 4.57 & {\footnotesize [3.96, 5.31]} & 1.00 & 0.78 & {\footnotesize [0.75, 0.81]} \\
\quad MiniCPM & 0.794 & 426 & 1.20 &  & 0.64 & 0.17 &  \\
\midrule
Range, 32 technologies & & 1.4--426 & \multicolumn{2}{l}{1.20--6.87} & 0.64--1.00 & \multicolumn{2}{l}{0.17--0.85} \\
\bottomrule
\end{tabular*}
\par\smallskip
\begin{tablenotes}
Clean inference-demand sample (77 models). Brackets (median $w$ and $s$): 2.5 and 97.5 percentiles over the technology's bootstrap draws,
recomputing the sample statistic in every draw (one technology draw shared by all models); no bracket: published point
estimate; the share with $w>1$ under the reference has the interval [0.92, 1.00]. A2: IsoFLOP-argmin (Approach 2) path. $M^*(10^{24})$: compute-optimal tokens per parameter at $10^{24}$ FLOP in the technology's own convention, an
extrapolation for every row. The ex-ante inclusion rule and the bootstrap draws are described in the text. Model-free
rows combine the design's own IsoFLOP path with its model-free $\sigma^*$ (Online Appendix Table~\ref{tab:app-mf-sens}).
Non-embedding and OLMo-convention technologies are evaluated at each model's own count.
\end{tablenotes}
\begin{tablenotes}[Source]
\citet{hoffmann2022training}; \citet{besiroglu2024chinchilla}; \citet{li2025predictableb}; \citet{gadre2024language};
\citet{bhagia2024establishing}; \citet{muennighoff2023scaling}; \citet{grattafiori2024llama}; \citet{marin2026ladders};
\citet{bi2024deepseek}; \citet{hu2024minicpm}; authors' calculations.
\end{tablenotes}
\end{table}


% -----------------------------------------------------------------------------------------------------
\subsection{Parameter Conventions and Mixture-of-Experts Models}\label{app:wedge-conv}

Each technology is evaluated at the model's own parameter count in the technology's convention: the total count for
the Chinchilla, Gadre et al.\ and laboratory technologies, the non-embedding count, from each model's configuration file,
for Farseer, MiniCPM and the non-embedding Chinchilla refit, and OLMo's count, which excludes only the input embedding,
for the OLMo ladder. A non-embedding technology ignores the output head, which costs $2N_{\rm head}$ FLOPs per token in
training and serving alike ($N_{\rm head}$ is vocabulary size times width). Counting it in FLOPs turns the ratio of
lifetime to training cost into $w(1+N_{\rm head}/N_{\rm ne})$.

Table~\ref{tab:app-conventions}, panel~A, shows that the sample median moves little across conventions, because
embeddings are at most 28 percent of parameters in the clean sample. For the eight clean models below one billion
parameters, the convention matters under one sweep but not the other. Comparing like with like, the non-embedding count
lowers their median wedge by 8 percent under the Chinchilla form (9.13 against 9.96) and by 36 percent under the Farseer
form (6.85 against 10.67). The Farseer gap exceeds the effect of tokenizer differences on $D$: under the reference,
$\partial\ln w/\partial\ln D=0.43$, so a 20 percent difference in tokens per byte moves the wedge by about 8 percent. We
do not normalize $D$ by bytes.

Mixture-of-experts models are excluded from the clean sample. Panel~B bounds their wedges under the reference. A dense
technology evaluated at active parameters treats a mixture of experts as parameter-augmenting and overstates its wedge;
evaluated at total parameters, it ignores the loss advantage of sparsity \citep{clark2022unified,krajewski2024scaling,
abnar2025parameters}. The median of the 12 models is 5.14 at active and 2.07 at total parameters. The bounds are widest
for the largest models: DeepSeek-V3 has a wedge of 3.50 or 1.01, and Kimi K2 3.79 or 0.87. Neither bound is an estimate
from a mixture-of-experts scaling law.

% Online Appendix Table F4: parameter-count conventions and mixture-of-experts bounds
% (output/tables/ra2_wedge_conventions.csv; ra2_wedge_models.csv columns w_chin_q, wtotalN_chin_q). Writer appendix_D_F.
\begin{table}[tp]
\centering
\caption{Parameter-Count Conventions and Mixture-of-Experts Models}
\label{tab:app-conventions}
\footnotesize
\setlength{\tabcolsep}{2.6pt}
\begin{tabular*}{\textwidth}{@{\extracolsep{\fill}}llrrrrr@{}}
\toprule
\multicolumn{7}{@{}l}{\textit{Panel A. Clean sample (77 models) under technologies in different parameter conventions}} \\[1pt]
 & & Median & Median $w$, & Share & Median & Median $w$, \\
Technology & Convention of $N$ & $w$ & head FLOPs & $w>1$ & $s$ & $N<1$B ($n=8$) \\
\midrule
Chinchilla, $\kappa$ free (reference) & Total & 3.99 &  & 0.97 & 0.75 & 15.10 \\
Chinchilla, $\kappa=1$ & Total & 3.34 &  & 1.00 & 0.70 & 9.96 \\
Farseer, $\kappa=1$ & Total & 2.64 &  & 0.88 & 0.62 & 10.67 \\
Chinchilla, $\kappa=1$ (refit) & Non-embedding & 3.65 & 3.74 & 1.00 & 0.73 & 9.13 \\
Farseer, $\kappa=1$ & Non-embedding & 2.74 & 2.86 & 1.00 & 0.63 & 6.85 \\
Farseer, $\kappa$ free & Non-embedding & 3.91 & 4.15 & 1.00 & 0.74 & 14.26 \\
MiniCPM (published law) & Non-embedding & 1.20 & 1.27 & 0.64 & 0.17 & 3.04 \\
\end{tabular*}
\begin{tabular*}{\textwidth}{@{\extracolsep{\fill}}lrrrrrrr@{}}
\midrule
\multicolumn{8}{@{}l}{\textit{Panel B. Mixture-of-experts models (excluded from the clean sample), reference technology}} \\[1pt]
 & $N$ (B), & $D$ & $M$, & \multicolumn{2}{c}{$w$} & \multicolumn{2}{c}{$s$} \\
\cmidrule(lr){5-6}\cmidrule(l){7-8}
Model (year) & active / total & (T) & active / total & Active $N$ & Total $N$ & Active $N$ & Total $N$ \\
\midrule
OLMoE 1B-7B (2024) & 1.3 / 7 & 5.1 & 3,946 / 741 & 9.18 & 4.51 & 0.89 & 0.78 \\
DeepSeek-V2-Lite (2024) & 2.4 / 16 & 5.7 & 2,375 / 363 & 7.40 & 3.34 & 0.86 & 0.70 \\
DeepSeekMoE 16B (2024) & 2.8 / 16 & 2.0 & 714 / 122 & 4.42 & 2.09 & 0.77 & 0.52 \\
Phi-3.5-MoE (2024) & 6.6 / 42 & 4.9 & 742 / 117 & 4.52 & 2.06 & 0.78 & 0.51 \\
DeepSeek-V2 (2024) & 21.0 / 236 & 8.1 & 386 / 34 & 3.43 & 1.23 & 0.71 & 0.19 \\
DeepSeek-V3 (2024) & 37.0 / 685 & 14.8 & 400 / 22 & 3.50 & 1.01 & 0.71 & 0.01 \\
Qwen3 30B-A3B (2025) & 3.3 / 31 & 36.0 & 10,909 / 1,179 & 14.30 & 5.56 & 0.93 & 0.82 \\
Llama 4 Scout (2025) & 17.0 / 109 & 40.0 & 2,353 / 368 & 7.46 & 3.40 & 0.87 & 0.71 \\
Qwen3 235B-A22B (2025) & 22.0 / 235 & 36.0 & 1,636 / 153 & 6.39 & 2.34 & 0.84 & 0.57 \\
GLM-4.5 (2025) & 32.0 / 358 & 23.0 & 719 / 64 & 4.50 & 1.61 & 0.78 & 0.38 \\
Llama 4 Maverick (2025) & 17.0 / 402 & 22.0 & 1,294 / 55 & 5.77 & 1.51 & 0.83 & 0.34 \\
Kimi K2 (2025) & 32.0 / 1,026 & 15.5 & 484 / 15 & 3.79 & 0.87 & 0.74 & $-$0.15 \\
\addlinespace
Median (12 models) & & & & 5.14 & 2.07 & 0.80 & 0.52 \\
\bottomrule
\end{tabular*}
\par\smallskip
\begin{tablenotes}
Panel A: medians over the clean inference-demand sample, each technology evaluated at the model's own count in the
technology's convention (non-embedding counts from each model's configuration file). With the output head counted in
FLOPs, the lifetime-to-training ratio for a non-embedding technology is $w(1+N_{\rm head}/N_{\rm ne})$. Last column:
the eight clean models below one billion parameters; comparing like with like, the non-embedding count lowers their
median wedge by 8 percent under the Chinchilla form (9.13 against 9.96) and by 36 percent under the Farseer form (6.85
against 10.67). Panel B: $w$ and $s=(w-1)/w$ under the reference technology (Chinchilla, $\kappa$ free), with $N$ equal
to active or to total parameters. A dense technology evaluated at active parameters treats a mixture-of-experts model
as parameter-augmenting and overstates its wedge; at total parameters it ignores the loss advantage of sparsity. The
two columns bound the wedge under a dense technology; neither is an estimate from a mixture-of-experts scaling law.
\end{tablenotes}
\begin{tablenotes}[Source]
Model cards, technical reports and Hugging Face configuration files; authors' calculations.
\end{tablenotes}
\end{table}


% -----------------------------------------------------------------------------------------------------
\subsection{The Developers' Own Technologies}\label{app:wedge-labown}

Where a developer has published scaling evidence, its own technology is primary for its models. For Meta this is its
IsoFLOP path, fitted to all ten budgets of Meta's figure, which reproduces the published law $D^*=0.299\,C^{0.537}$
\citep{grattafiori2024llama}, combined with the model-free curvature of Meta's own IsoFLOP profiles, 0.660 (standard
error 0.023). For AI2 it is the OLMo ladder with $\kappa$ free \citep{bhagia2024establishing}; for DeepSeek, its
published law with the reference curvature, since DeepSeek published no curvature; and for Marin, the Nemotron-CC ladder
path with its model-free curvature (Marin 8B's model card lists Nemotron-CC and DCLM among its data). These technologies
cover 22 models: Meta 10, AI2 9, DeepSeek 2 and Marin 1. For 11 of them the developer's law was published before the
model was planned. For the other 11 (Llama~1 and~2, OLMo~1 and~1.7, Marin 8B) it came later, so it describes the
developer's technology but not the law it planned with.

Table~\ref{tab:app-labown}, panel~A, reports two model-free estimators of the designs' curvature. The primary one
(Section~\ref{sec:tech}) fits local quadratics near the path on budgets whose minimum is bracketed by two runs on each
side; the second fits one quadratic to all runs of every budget, including unbracketed ones, and pools a ratio of sums.
They agree on Llama~3 (0.660 and 0.653) and differ by 0.05 to 0.08 on Marin; after the second estimator's
simulation-based bias correction (to 0.66--0.68 on every design), the Marin gap is 0.03 to 0.05, within the range of the
primary estimator's own specifications (0.64 to 0.72). We use the primary estimator, whose Monte Carlo bias is
negligible (Table~\ref{tab:app-mf-sens}).

On the 22 models, the median share is 0.61 under the developers' technologies and 0.65 under the reference (0.61 and 0.67
on the 11 contemporaneous ones). Panel~B shows selected models; four points qualify them.
\begin{itemize}
\item[] \emph{Meta.} Two of Meta's ten budgets ($3\times10^{21}$ and $10^{22}$ FLOP, six runs each) do not bracket their
minimum. On the eight that do, the path slope is 0.501 rather than 0.463 and $M^*(10^{24})$ is 15 rather than 31; Llama~3
8B's wedge becomes 12.2 and the 405B's 1.64. We keep the ten-budget path because the developer's technology is meant to
be the law Meta planned with, which sized the 405B model; its wedge of 0.97 is on Meta's path by construction. With the
second curvature estimator after its simulation-based correction ($\sigma^*=0.678$), Llama~3 8B's wedge is 7.07.
\item[] \emph{AI2.} The $\kappa$-free ladder has a steep curvature ($1/\sigma^*-1=0.84$), which makes the wedges of small
models extreme: 20.2 for OLMo~2 1B, against 2.0 under the ladder with $\kappa=1$. OLMo~2 32B's interval, [0.96, 1.72],
contains one.
\item[] \emph{DeepSeek.} The published law gives 3.42 for the 7B and 1.34 for the 67B model; the intervals reflect only
the reference curvature.
\item[] \emph{Marin.} The curvature estimator matters: Marin 8B's wedge is 5.08 with the primary estimate and 6.66 with
the corrected second one, and 1.43 to 6.66 across Marin's own alternatives.
\end{itemize}

% Online Appendix Table F5: lab-own technologies (output/tables/ra2_wedge_modelfree_sigma.csv, ra2_wedge_pi_anchors.csv,
% ra2_wedge_labown.csv; ra1_modelfree_isoflop_summary.csv for the sigma* used). Writer appendix_D_F, 2026-09-24.
\begin{table}[tp]
\centering
\caption{Lab-Own Technologies: IsoFLOP Curvature, Paths and Wedges}
\label{tab:app-labown}
\footnotesize
\setlength{\tabcolsep}{1.6pt}
\begin{tabular*}{\textwidth}{@{\extracolsep{\fill}}lcccccc@{}}
\toprule
\multicolumn{7}{@{}l}{\textit{Panel A. Model-free curvature and compute-optimal paths of the IsoFLOP designs}} \\[1pt]
 & Budgets & Largest & $\sigma^*$ used & $\sigma^*$, global & Path & $M^*$ at largest \\
Design & (valid) & budget & (s.e.) & quadratics [95\%] & slope $a$ & budget [95\%] \\
\midrule
Llama 3 (Meta) & 10 (8) & $10^{22}$ & 0.660 (0.023) & 0.653 [0.639, 0.666] & 0.463 & 22.3 [20.8, 23.7] \\
Marin, Comma & 7 (5) & $3\times10^{20}$ & 0.700 (0.027) & 0.650 [0.611, 0.669] & 0.405 & 12.4 [10.7, 13.9] \\
Marin, DCLM & 7 (7) & $3\times10^{20}$ & 0.713 (0.027) & 0.629 [0.586, 0.653] & 0.424 & 9.4 [8.0, 10.6] \\
Marin, Nemotron-CC & 8 (7) & $3\times10^{20}$ & 0.705 (0.023) & 0.644 [0.613, 0.666] & 0.422 & 10.0 [8.7, 11.1] \\
Chinchilla profiles & 9 (9) & $3\times10^{21}$ & 0.673 (0.027) & 0.610 [0.588, 0.644] & 0.492 & 22.7 [18.2, 29.4] \\
\end{tabular*}
\begin{tabular*}{\textwidth}{@{\extracolsep{\fill}}llccccc@{}}
\midrule
\multicolumn{7}{@{}l}{\textit{Panel B. Wedges of selected models under their developer's own technology}} \\[1pt]
 & Lab law & Lab-own $w$ & Lab-own & Range over lab's & Reference & Reference \\
Model & vintage & [95\%] & $s$ & alternatives & $w$ & $s$ \\
\midrule
Llama 3 8B & contemp. & 8.36 [5.67, 12.76] & 0.88 & 3.09--12.16 & 6.73 & 0.85 \\
Llama 3 70B & contemp. & 2.51 [2.09, 3.02] & 0.60 & 1.63--3.96 & 2.67 & 0.63 \\
Llama 3.1 405B & contemp. & 0.97 [0.88, 1.08] & $-$0.03 & 0.97--1.64 & 1.29 & 0.23 \\
Llama 2 70B & ex post & 0.97 [0.91, 1.04] & $-$0.03 & 0.97--1.42 & 1.13 & 0.12 \\
OLMo 2 1B & contemp. & 20.2 [14.5, 35.2] & 0.95 & 2.01--20.2 & 7.83 & 0.87 \\
OLMo 2 7B & contemp. & 4.05 [3.22, 5.32] & 0.75 & 1.28--4.05 & 3.99 & 0.75 \\
OLMo 2 32B & contemp. & 1.28 [0.96, 1.72] & 0.22 & 0.91--1.28 & 2.54 & 0.61 \\
DeepSeek LLM 7B & contemp. & 3.42 [3.05, 3.87] & 0.71 & & 2.99 & 0.67 \\
DeepSeek LLM 67B & contemp. & 1.34 [1.31, 1.38] & 0.26 & & 1.15 & 0.13 \\
Marin 8B & ex post & 5.08 [3.10, 8.71] & 0.80 & 1.43--6.66 & 6.26 & 0.84 \\
\addlinespace
Median, 22 models & & & 0.61 & & & 0.65 \\
Median, 11 contemp. & & & 0.61 & & & 0.67 \\
\bottomrule
\end{tabular*}
\par\smallskip
\begin{tablenotes}
Panel A: $\sigma^*$ used is the model-free random-effects estimate of Section~\ref{sec:tech} over budgets whose minimum
is bracketed by two runs on each side (Online Appendix Table~\ref{tab:app-mf-sens}). Global quadratics: a second
estimator that fits one quadratic to all runs of each budget, including unbracketed ones, and pools a ratio of sums;
design-conditional wild bootstrap, 999 draws. Its simulation-based bias correction gives 0.662--0.678, and it is used
only in sensitivity rows; for Chinchilla it uses the 123-point snapped subset of the IsoFLOP compilation of \citet{czech2026problems} (eight budgets) as a cross-check. Path slope and $M^*$: Approach-2 argmins regressed on $\ln C$ over all budgets; brackets:
wild bootstrap. Meta's path reproduces the published law $D^*=0.299\,C^{0.537}$; fitted on its eight bracketed budgets
it has slope 0.501 and $M^*(10^{24})=15$ instead of 31. Panel B: lab-own technologies are Meta's IsoFLOP path with the
model-free curvature; AI2's OLMo ladder with $\kappa$ free; DeepSeek LLM's published allocation law in its own units
with the reference curvature (no curvature published); and Marin's Nemotron-CC path with the model-free curvature.
Vintage: whether the lab's law was published before the model was planned (contemporaneous) or afterward (ex post).
Llama 3.1 405B lies on Meta's path by construction, because Meta sized it with this law. Alternatives: Meta's $\kappa=1$
law, primal and path on eight budgets; AI2's $\kappa=1$ ladder; Marin's other corpora and $\kappa=1$ rows. Reference:
Chinchilla with $\kappa$ free. Medians are over the 22 models with a lab-own technology (Meta 10, AI2 9, DeepSeek 2,
Marin 1).
\end{tablenotes}
\begin{tablenotes}[Source]
\citet{grattafiori2024llama}; \citet{czech2026llama3isoflop}; \citet{marin2026ladders}; \citet{bhagia2024establishing};
\citet{bi2024deepseek}; authors' calculations.
\end{tablenotes}
\end{table}


% -----------------------------------------------------------------------------------------------------
\subsection{Extrapolation and Partial Identification}\label{app:wedge-pi}

\emph{Extrapolation in $M$.}---Farseer's near-factorial grid reaches $M=2{,}570$ in non-embedding parameters, which
allows a direct check of parametric wedges at high ratios. Table~\ref{tab:app-farseer-w} compares the local wedge of
Proposition~\ref{prop:modelfree}(i), from a kernel-weighted local quadratic of the loss surface, with the wedges of
parametric forms fitted to the same runs. Every form understates the wedge at high $M$, and the gap grows with $M$. At
$M\ge1{,}024$, where the local wedge is 7.5, the $\kappa$-free form's wedge is 0.70 times the local one and the Chinchilla
form's 0.44 times; the Chinchilla form fitted only to runs with $M\le100$ gives 0.26 times. The sign of the gap survives
every bandwidth, but its size does not: for the $\kappa$-free form, $\ln(w_{\rm param}/w_{\rm local})$ runs from $-0.25$ at
the narrower cross-validated bandwidth to $-0.67$ at twice the primary one (panel~B). The source of the gap is curvature
of the wedge itself. In a regression of the local $\ln w$ on $u=\ln(M/M^*(C))$ and $u^2$, the quadratic term is 0.020
(standard error 0.001) and positive at every bandwidth, so $\ln w$ is convex in $u$, whereas every member of the $\kappa$
family makes it linear.

Applying Farseer's gaps by bin of $M$ to the 57 clean models inside Farseer's range of $M$ raises their median wedge
under the $\kappa$-free Farseer technology from 3.10 to 3.56, and their median share from 0.68 to 0.72; at the narrower
bandwidth, to 3.49 and 0.71. Under the Chinchilla form fitted to Farseer the median wedge rises from 2.29 to 3.73. The
correction is negative for the 56 percent of these models with $M\ge256$ and slightly positive, about $+0.05$ in logs,
at $M$ of 16 to 64. This transfers one recipe's gaps to other developers' models and to larger $N$ and $C$, so it signs
the error rather than correcting levels. For the 20 clean models beyond Farseer's range of $M$ there is no model-free
evidence, and revealed preference recovers the technology a developer believed, which may be the parametric one
(Remark~\ref{rem:A-pi-param}).

% Online Appendix table: model-free local wedge against parametric wedges on Farseer (condensed from
% output/tables/ra1_modelfree_farseer_w.tex; values from ra1_modelfree_farseer_delta.csv and
% ra1_modelfree_farseer_sensitivity.csv). Writer appendix_D_F, 2026-09-24.
\begin{table}[tp]
\centering
\caption{Is the Wedge Extrapolated Correctly? Local and Parametric Wedges on Farseer}
\label{tab:app-farseer-w}
\footnotesize
\setlength{\tabcolsep}{3pt}
\begin{tabular*}{\textwidth}{@{\extracolsep{\fill}}lcccccc@{}}
\toprule
 & Local & \multicolumn{5}{c}{$\ln(w_{\rm param}/w_{\rm local})$} \\
\cmidrule(l){3-7}
 & $\ln w$ & \multicolumn{2}{c}{Chinchilla form} & \multicolumn{2}{c}{$\kappa$ free} & Farseer \\
\cmidrule(lr){3-4}\cmidrule(lr){5-6}
$M=D/N$ & ($w$) & All runs & $M\le100$ & All runs & $M\le100$ & Eq.~3 \\
\midrule
\multicolumn{7}{@{}l}{\textit{Panel A. By bin of $M$; non-embedding parameters, primary bandwidth}} \\
$M<16$ & $-$0.69 (0.50) & 0.16 & 0.08 & $-$0.02 & 0.02 & $-$0.03 \\
16--64 & 0.06 (1.06) & 0.01 & $-$0.23 & 0.05 & $-$0.04 & 0.01 \\
64--256 & 0.68 (1.96) & $-$0.19 & $-$0.55 & 0.00 & $-$0.18 & $-$0.01 \\
256--1,024 & 1.37 (3.93) & $-$0.49 & $-$0.95 & $-$0.14 & $-$0.42 & $-$0.13 \\
$\ge$1,024 & 2.01 (7.48) & $-$0.82 & $-$1.36 & $-$0.36 & $-$0.70 & $-$0.31 \\
\addlinespace
\multicolumn{7}{@{}l}{\textit{Panel B. The $M\ge1{,}024$ bin under alternative bandwidths and conventions}} \\
Extended-grid optimal $h_N$ & 1.84 & $-$0.70 & $-$1.22 & $-$0.25 & $-$0.58 & $-$0.20 \\
$h_N\times0.75$ & 1.96 & $-$0.78 & $-$1.32 & $-$0.32 & $-$0.66 & $-$0.27 \\
Bandwidths $\times1.5$ & 2.16 & $-$0.96 & $-$1.50 & $-$0.49 & $-$0.84 & $-$0.44 \\
Bandwidths $\times2$ & 2.35 & $-$1.14 & $-$1.68 & $-$0.67 & $-$1.02 & $-$0.63 \\
With embeddings in $N$ & 2.74 & $-$1.17 & $-$1.67 & $-$0.70 & $-$1.00 & $-$0.62 \\
\bottomrule
\end{tabular*}
\par\smallskip
\begin{tablenotes}
Farseer's 404 released runs \citep{li2025predictableb}. Local $\ln w$: Gaussian product-kernel local quadratic of $\ln L$
in $(\ln N,\ln D)$, with $w=\partial_{\ln N}\ln L/\partial_{\ln D}\ln L$ (Proposition~\ref{prop:modelfree}(i); invariant
to monotone transformations of the loss), averaged over grid points inside each model size's observed token range with
a kernel effective sample size of at least 15; primary bandwidths $h_N=0.24$ and $h_D=0.74$ in logs, from
leave-one-out cross-validation over multiples of the standard deviation ($h_N$ at the smallest multiple on that grid;
the optimum on an extended grid is $h_N=0.148$). Other columns: log ratio of the wedge implied by a parametric form
fitted to the same runs to the local wedge at the same points; negative values mean that the form understates the
wedge. Chinchilla form and $\kappa$ family by Huber loss, Farseer's Eq.~3 by nonlinear least squares; ``$M\le100$'':
fitted only to runs with $M\le100$ and extrapolated. In the $M\ge1{,}024$ row, the 95 percent intervals are [$-$0.90, $-$0.75], [$-$1.44, $-$1.28], [$-$0.43, $-$0.28], [$-$0.79, $-$0.64] and [$-$0.38, $-$0.23] in column order, and those of the other rows are narrower: basic intervals from a wild cluster bootstrap
by model size (Webb weights, 999 draws, every model refitted); they exclude the smoothing bias of the local estimator,
which Panel B shows instead. Only the models of 0.10 to 0.34 billion parameters reach $M\ge1{,}024$.
\end{tablenotes}
\begin{tablenotes}[Source]
\citet{li2025predictableb}; authors' calculations.
\end{tablenotes}
\end{table}


\emph{Extrapolation in compute.}---No design observes $M^*(C)$ above $10^{22}$ FLOP; the largest IsoFLOP budgets are
$3\times10^{21}$ (Chinchilla), $10^{22}$ (Llama~3) and $3\times10^{20}$ FLOP (Marin; DeepSeek LLM). The 402 billion
parameters and 16.55 trillion tokens that Meta reports at $3.8\times10^{25}$ FLOP extrapolate its law and are not used
as an observation. We anchor $M^*$ at each design's largest budget, with a 95 percent wild-bootstrap interval: 22.7
[18.2, 29.4] at $3\times10^{21}$ FLOP for Chinchilla's profiles, 22.3 [20.8, 23.7] at $10^{22}$ for Llama~3, 12.4, 9.4
and 10.0 at $3\times10^{20}$ for Marin's three corpora, and DeepSeek's published law. Proposition~\ref{prop:pi} then gives
an identified set for $M^*(C)$ once the elasticity of $M^*$ with respect to compute, $e=1-2a$, is bounded.
Table~\ref{tab:app-pi} reports four sets of assumptions. PI-1 lets $e$ lie anywhere in the range spanned by the IsoFLOP
paths and the published laws, including DeepSeek's three data sets and MiniCPM, $[-0.16,0.19]$. PI-2 adds that $M^*$ does
not fall with compute. PI-3 uses a developer's own anchor where one exists. PI-4 admits every ex-ante technology's
in-support $M^*$ as an anchor, with the widest range of slopes, $[-0.16,0.31]$.

By Proposition~\ref{prop:pi}(iii), the sign of $w-1$ is the sign of $\ln M-\ln M^*(C)$ for any technology with
single-peaked IsoFLOP profiles, so over-training is identified whenever $M$ exceeds the upper bound of the set. That is
the case for 86 percent of clean models under PI-1 and PI-2, 77 percent under PI-3 and 23 percent under PI-4. No model's
wedge is identified as below one under any set. The 11 models whose sign is not identified under PI-1 all have $M\le96$:
Llama~2 70B, LLaMA 30B and 65B, Llama~3.1 405B, Qwen-72B, Qwen2-72B, Yi-34B, Falcon 40B and 180B, MPT-30B and DeepSeek
LLM 67B. Under PI-3, AI2's models join them, because the OLMo ladder's own anchor is imprecise. Magnitudes need a
curvature; with $1/\sigma^*-1$ between 0.40 and 0.52, the range of the $\kappa$-free and model-free estimates, the medians
of the models' lower and upper bounds on $s$ are 0.59 and 0.93 under PI-1.

Panel~C repeats the logic on the 173 verified models in their own units, with anchors fitted by least squares to the
IsoFLOP minima and wild bootstrap-$t$ intervals (Online Appendix~\ref{app:proofs}, Remark~\ref{rem:A-pi-illus}). It
illustrates the bounds rather than estimating them, because units are not harmonized. With a cross-sweep range of slopes,
over-training is identified for 77 to 81 percent of the models beyond each anchor's largest budget; with normal inputs
alone ($|e|\le1$), for about a third. The two anchors' point extrapolations disagree for Llama~3.1 405B, which is
over-trained against the Chinchilla anchor and under-trained against Meta's, and its sign is not identified under any
set-valued assumption. The two Chinchilla anchors differ, 22.7 in Panel~A and 20.5 [11.4, 36.8] in Panel~C, because
Panel~A uses the eight-budget compilation of \citet{czech2026problems} and a percentile interval, and Panel~C all nine
budgets of the digitization and a bootstrap-$t$ interval with HC2-rescaled residuals.

% Online Appendix Table F6: partial identification of M*(C) at frontier scale (output/tables/ra2_wedge_pi.tex,
% ra2_wedge_pi_*.csv; Panel C from ra5_theory_pi_sign.csv/.tex). Writer appendix_D_F, 2026-09-24.
\begin{table}[tp]
\centering
\caption{Partial Identification of the Compute-Optimal Ratio at Frontier Scale}
\label{tab:app-pi}
\footnotesize
\setlength{\tabcolsep}{2pt}
\begin{tabular*}{\textwidth}{@{\extracolsep{\fill}}lcccccc@{}}
\toprule
\multicolumn{7}{@{}l}{\textit{Panel A. Bounds on $M^*(C)$ at frontier compute (harmonized units; total-$N$ anchors)}} \\[1pt]
Assumption set & $e$ range & $M^*(10^{23})$ & $M^*(10^{24})$ & $M^*(10^{25})$ & & Anchors \\
\midrule
PI-1: IsoFLOP anchors & [$-$0.16, 0.19] & [3.2, 57.3] & [2.3, 88.8] & [1.6, 138] & & 5 \\
PI-2: PI-1, nondecreasing & [0.00, 0.19] & [8.0, 57.3] & [8.0, 88.8] & [8.0, 138] & & 5 \\
PI-3: Meta's own anchor & [$-$0.16, 0.19] & [14.5, 36.8] & [10.1, 57.0] & [7.1, 88.3] & & 1 \\
PI-4: all technologies & [$-$0.16, 0.31] & [0.4, 1,609] & [0.3, 3,293] & [0.2, 6,739] & & 30 \\
\midrule
\multicolumn{7}{@{}l}{\textit{Panel B. Sign of $w-1$ on the clean sample ($n=77$), each anchor in its own convention}} \\[1pt]
 & & $w>1$ & Not & \multicolumn{2}{c}{Median bound on $s$} & \\
\cmidrule(lr){5-6}
Assumption set & $e$ range & identified & identified & Lower & Upper & \\
\midrule
PI-1 & [$-$0.16, 0.19] & 0.86 & 0.14 & 0.59 & 0.93 & \\
PI-2 & [0.00, 0.19] & 0.86 & 0.14 & 0.59 & 0.89 & \\
PI-3 & [$-$0.16, 0.19] & 0.77 & 0.23 & 0.55 & 0.94 & \\
PI-4 & [$-$0.16, 0.31] & 0.23 & 0.77 & $-$1.43 & 0.98 & \\
\midrule
\multicolumn{7}{@{}l}{\textit{Panel C. Illustration in unharmonized units: verified models beyond the anchor's largest budget}} \\[1pt]
 & & & \multicolumn{3}{c}{Share of models} & \\
\cmidrule(lr){4-6}
Assumption & $e$ range & $M^*(10^{24})$ & $w>1$ & $w<1$ & Not identified & \\
\midrule
\multicolumn{7}{@{}l}{Chinchilla anchor ($n=149$)} \\
\quad Normal inputs & [$-$1.00, 1.00] & [0.03, 12,651] & 0.31 & 0.00 & 0.69 & \\
\quad Cross-sweep slopes & [$-$0.14, 0.26] & [5.0, 168] & 0.77 & 0.01 & 0.21 & \\
\quad Anchor's own slope & [$-$0.16, 0.17] & [4.3, 100] & 0.83 & 0.01 & 0.16 & \\
\quad Point extrapolation & 0.005 & 21.1 & 0.97 & 0.03 & 0.00 & \\
\multicolumn{7}{@{}l}{Llama~3 anchor ($n=137$)} \\
\quad Normal inputs & [$-$1.00, 1.00] & [0.12, 4,020] & 0.34 & 0.00 & 0.66 & \\
\quad Cross-sweep slopes & [$-$0.14, 0.26] & [6.4, 133] & 0.81 & 0.01 & 0.18 & \\
\quad Anchor's own slope & [$-$0.09, 0.24] & [7.9, 119] & 0.84 & 0.02 & 0.14 & \\
\quad Point extrapolation & 0.074 & 31.3 & 0.90 & 0.10 & 0.00 & \\
\bottomrule
\end{tabular*}
\par\smallskip
\begin{tablenotes}
$M^*(C)$ is never observed beyond $10^{22}$ FLOP: the largest IsoFLOP budgets are $3\times10^{21}$ (Chinchilla),
$10^{22}$ (Llama 3) and $3\times10^{20}$ (Marin; DeepSeek LLM). Anchors: each design's Approach-2 path at its largest
budget with its 95 percent wild-bootstrap interval (DeepSeek: published law). Extrapolation (Proposition~\ref{prop:pi}):
$\ln M^*(C)=\ln M^*(C_0)+e\ln(C/C_0)$ with the path elasticity $e=1-2a$ anywhere in the stated range. PI-1: range of
$e$ over the IsoFLOP paths and published laws (including DeepSeek LLM's three data sets and MiniCPM); PI-2 adds $M^*$
nondecreasing in $C$; PI-3 uses the model's own lab's anchor where it exists (Meta, Marin, DeepSeek, AI2) and PI-1
otherwise; PI-4 uses every ex-ante technology's in-support $M^*$ and all path slopes. Llama 3's 402B parameters and
16.55T tokens at $3.8\times10^{25}$ FLOP extrapolate Meta's law and are not used as an anchor. By
Proposition~\ref{prop:pi}(iii), $w>1$ is identified when $\ln M$ exceeds the upper bound for every anchor, for any
technology with single-peaked IsoFLOP profiles; no model's $w<1$ is identified in Panel B. Bounds on $s$ use
$1/\sigma^*-1\in[0.40,0.52]$, the range of the $\kappa$-free and model-free estimates, and are medians of the models'
bounds. Panel C repeats the logic on the verified sample with each model's own parameter and token conventions and
anchors fitted to the IsoFLOP minima (wild bootstrap-$t$ intervals with HC2-rescaled residuals, $M^*=20.5$ [11.4, 36.8]
at $2.9\times10^{21}$ for Chinchilla and 22.3 [12.2, 40.2] at $10^{22}$ for Llama 3); cross-sweep slopes: allocation
exponents 0.37--0.57; anchor's own slope: its 95 percent interval. It illustrates the bounds, not the paper's estimate.
\end{tablenotes}
\begin{tablenotes}[Source]
\citet{hoffmann2022training}; \citet{grattafiori2024llama}; \citet{marin2026ladders}; \citet{bi2024deepseek};
\citet{hu2024minicpm}; authors' calculations.
\end{tablenotes}
\end{table}


% -----------------------------------------------------------------------------------------------------
\subsection{Conduct}\label{app:wedge-conduct}

Table~\ref{tab:app-conduct} reports the three conduct tests of Section~\ref{sec:wedge-conduct}. All cluster by parent
developer and use a restricted wild cluster bootstrap-$t$ with the six-point weights of \citet{webb2023reworking}
\citep{cameron2008bootstrap}. Because $\ln\hat w$ is a function of $(N,C)$ under one technology, every test concerns the
choice of size at given compute.

\emph{Open and closed models.}---The test uses the production-scale universe of Online Appendix~\ref{app:data}
($6ND\ge10^{21}$ FLOP, 2019--2026), restricted to models with confident parameter and token counts and without
mixture-of-experts or non-transformer architectures, with the clean-sample exclusions applied to verified models: 186
models from 67 developers, 23 of which released at least one closed model. At given compute and year, open-weight models
have log wedges 0.52 higher ($p=0.011$), and 0.71 higher among releases since 2023 ($p=0.001$). On the share the premium is
0.29 ($p=0.097$). Closed developers have disclosed few allocations since 2023: six in 2024 and two in 2025.

\emph{Serving footprint and deployment target.}---A developer counts as serving if, at the model's release, it operated a
commercial interface or consumer product that served its own language models. Meta is coded as serving from 27 September
2023, the launch of Meta AI, so Llama~1 and~2 are not; Alibaba, Google, Microsoft, IBM, MosaicML, DeepSeek (from 2024),
01.AI (from 2024), Moonshot, Zhipu and Together are coded as serving; TII, Hugging Face, H2O.ai, AI2, Stability AI, Swiss
AI, Marin, M-A-P, TinyLlama and EleutherAI are not. Two dates are judgment calls: phi-1.5, released under a research
license in September 2023, is coded as served, and 01.AI's interface is dated to 2024. A model has an on-device target if
its card names on-device, mobile or phone use, and a local target if it names laptops or desktops. In the clean sample,
33 models come from 7 serving developers; 9 models have an on-device and 3 a local target, but only two and three
developers carry these attributes. With three or fewer treated clusters the wild cluster bootstrap is unreliable
\citep{mackinnon2018wild}, so only the serving test is informative, and with 18 clusters it has little power: 0.45
(standard error 0.26; $p=0.39$). Adding the research suites and replications raises the sample to 136 models and 26
developers and gives 0.47 ($p=0.070$), and 0.14 on the share ($p=0.046$). That sample adds suites whose developers do not
serve, however, so the contrast is partly between suites and products.

\emph{Tiers.}---Windows around quantized memory boundaries (2.4--3.3, 6.5--9.5, 11.5--14.9, 26--32.9 and 65--72.9
billion parameters) hold 51 percent of the 220 open-weight production-scale models released in 2023--2026, against 13
percent under a smooth polynomial counterfactual \citep{chetty2011adjustment,kleven2016bunching}; the excess mass is 83
models (standard error 9.8), 45 of them at 6.5--9.5 billion. Position relative to the tiers does not predict the wedge
beyond what the mechanical relation between $N$ and $\ln w$ implies. Models inside a window have lower wedges given
compute ($-0.39$, standard error 0.28), and 99.7 percent of 999 randomly shifted tier menus give a larger coefficient, so
the real windows are, if anything, less over-trained than placebo windows. The log distance to the next cap has a
coefficient of 0.46, against a placebo mean of 0.45.
Tier menus shape the choice of $N$; the tier model's prediction for the wedge is not visible at given compute.

% Online Appendix Table E6 (round 3, WP4b, fix list W14, W23(l), W24(h)): copy of output/tables/rb2_decisions_conduct.tex
% (49 budget-level decision units; rb2_decisions re-run, reviewed) with Panel D added from
% output/tables/rb5_units_serving_ri.csv and rb5_units_serving_ri_dev.csv (rb5_units, reviewed).
\begin{table}[tp]
\centering
\caption{Conduct Tests with Model-Level Serving and Decision Units}
\label{tab:app-conduct}
\footnotesize
\setlength{\tabcolsep}{3pt}
\begin{tabular}{@{}lrrrrll@{}}
\toprule
& $n$ & Clusters & Treated & Coef. & (s.e.) & WCR $p$ \\
\midrule
\multicolumn{7}{@{}l}{\textit{Panel A. Serving, outcome $\ln\hat w$ given $\ln C$ and year effects}} \\[1pt]
Models: developer-level code & 77 & 18 & 7 & 0.43 & (0.27) & 0.481 \\
Models: model-level code & 77 & 18 & 6 & 0.56 & (0.24) & 0.413 \\
Models: model-level, ambiguous sizes served & 77 & 18 & 7 & 0.62 & (0.24) & 0.220 \\
Models without Alibaba & 58 & 17 & 5 & $-$0.09 & (0.28) & 0.751 \\
Decisions: any member served & 49 & 18 & 6 & 0.28 & (0.29) & 0.673 \\
Decisions: all members served & 49 & 18 & 4 & 0.49 & (0.17) & 0.159 \\
Decisions without Alibaba & 43 & 17 & 5 & $-$0.20 & (0.30) & 0.572 \\
\addlinespace
\multicolumn{7}{@{}l}{\textit{Panel B. Open-weight premium, production-scale universe, outcome $\ln\hat w$}} \\[1pt]
Models & 186 & 67 & 55 & 0.51 & (0.18) & 0.012 \\
Decisions & 159 & 67 & 55 & 0.45 & (0.17) & 0.018 \\
Decisions, 2023 and later & 123 & 60 & 52 & 0.63 & (0.15) & 0.001 \\
\addlinespace
\multicolumn{7}{@{}l}{\textit{Panel C. Descriptive: median $s$ (reference)}} \\[1pt]
Served (model-level) & 26 & 6 & & 0.83 & {\scriptsize[0.63, 0.91]} & \\
Not served & 51 & 15 & & 0.74 & {\scriptsize[0.60, 0.84]} & \\
Served, without Alibaba & 7 & 5 & & 0.55 & {\scriptsize[0.21, 0.78]} & \\
Served, in a tier window & 18 & 5 & & 0.74 & {\scriptsize[0.57, 0.86]} & \\
On-device target & 9 & 2 & & 0.90 & {\scriptsize[0.84, 0.92]} & \\
Server or unspecified & 65 & 16 & & 0.70 & {\scriptsize[0.57, 0.85]} & \\
Decisions, served & 11 & 6 & & 0.68 & {\scriptsize[0.30, 0.79]} & \\
Decisions, not served & 38 & 14 & & 0.74 & {\scriptsize[0.61, 0.84]} & \\
\addlinespace
\multicolumn{7}{@{}l}{\textit{Panel D. Randomization inference for the serving coefficient of Panel A (two-sided)}} \\[1pt]
& $n$ & Strata & Treated & Coef. & RI $p$ & RI-$t$ $p$ \\
Models, strata: year $\times$ compute tercile & 77 & 9 & 6 & 0.56 & $<$0.001 & 0.06 \\
Models, strata: year $\times$ compute half & 77 & 6 & 6 & 0.56 & $<$0.001 & 0.07 \\
Models, strata: year & 77 & 3 & 6 & 0.56 & $<$0.001 & 0.07 \\
Models without Alibaba, year $\times$ tercile & 58 & 9 & 5 & $-$0.09 & 0.63 & 0.76 \\
Decisions, year $\times$ tercile & 49 & 9 & 6 & 0.28 & 0.17 & 0.38 \\
Decisions, strata: year & 49 & 3 & 6 & 0.28 & 0.11 & 0.40 \\
Developers, slope weighted by models & 77 & 18 & 6 & 0.56 & 0.10 & \\
Developers, unweighted slope & 77 & 18 & 6 & 0.24 & 0.28 & \\
Developers without Alibaba, weighted & 58 & 17 & 5 & 0.01 & 0.94 & \\
\bottomrule
\end{tabular}
\medskip
\begin{tablenotes}
Panels A--B: restricted wild cluster bootstrap-t by developer, Webb weights, $B=9{,}999$; CRV1 standard errors in parentheses; treated = developer clusters with the attribute (with three or fewer the $p$-value is not interpretable). Model-level serving code: this size served first-party (per-token API or consumer product) within 180 days of release; ambiguous: Meta AI ``built with Llama 3'' (8B, 70B) and 01.AI's early-access API (Yi-34B). Decisions: members of a family that share one token budget enter once (family share; total compute; first release year); in Panel B, the clean sample's common budgets enter once and other models of the production-scale universe are decisions of their own. Panel C: interquartile range in brackets; columns 2--3 give $n$ and developers. Panel D: randomization inference \citep{mackinnon2020randomization}, 9,999 permutations; RI $p$ for the coefficient, RI-$t$ $p$ for its CRV1 $t$-statistic. The rows for models and decisions permute the serving code across units within the strata shown (column 3: number of strata), which treats each size as a separate assignment although one developer decides for all its sizes; the developer rows permute the developers' shares of served models across the 18 developers (column 3: developers) and use the developer-level slope of the residuals of $\ln\hat w$ on $\ln C$ and year effects.
\end{tablenotes}
\end{table}


% -----------------------------------------------------------------------------------------------------
\subsection{Family Token Budgets}\label{app:wedge-family}

A family has a common token budget if its members' budgets lie within 10 percent of each other; Llama~3 and~3.1 count
as one herd. Eleven families of the clean sample, with 32 models, qualify; among all 173 verified models, 23 families
with 83 models do.

If sizes are chosen at the margin and $D$ is shared, the family's first-order condition in $D$ identifies one number, a
weighted average of the members' values of compactness (Proposition~\ref{prop:A-family}). Without data costs,
$\sum_i\pi_im_{N,i}=\bar w_H-1$, where $\bar w_H=[\sum_i\omega_i/w_i]^{-1}$ is the harmonic mean of the member wedges
weighted by training compute and $\pi_i\propto\omega_i/w_i$. Member wedges then differ only through the technology,
$\ln w_i-\ln w_j=-\alpha(\ln N_i-\ln N_j)$, and are not revealed preferences. If sizes come from a menu, the condition in
$D$ says nothing about the members' values of compactness. In three family problems solved by brute force, the condition
holds to within $2.2\times10^{-8}$.

Table~\ref{tab:app-family-budgets} reports the family shares $1-1/\bar w_H$. They run from 0.18 for DeepSeek LLM to 0.89
for Qwen3. Because the flagship carries most of a family's compute, the family share lies far below its members' typical
shares: the Llama~3 herd's members have wedges from 1.30 to 6.73, but the herd's weighted condition gives 1.42, a share
of 0.29. The family number equals the compute-weighted mean of the members' lifetime-to-training ratios when those ratios
are uncorrelated with $1/w_i$, and it is a lower bound when smaller members, which are more over-trained, are also
planned for more serving.

% Online Appendix Table F7: family-level token budgets (output/tables/ra2_wedge_family_level_W.csv). Writer appendix_D_F.
\begin{table}[tp]
\centering
\caption{Families with a Common Token Budget: What the Family Reveals}
\label{tab:app-family-budgets}
\footnotesize
\setlength{\tabcolsep}{3pt}
\begin{tabular*}{\textwidth}{@{\extracolsep{\fill}}lcccccc@{}}
\toprule
 & & Max./min. & Member wedges & Compute-weighted & Family & Family \\
Family & Members & $D$ & $w_i$ (range) & mean of $w_i$ & $\bar w_H$ & share $1-1/\bar w_H$ \\
\midrule
Qwen3 & 5 & 1.00 & 7.57--29.56 & 10.05 & 9.13 & 0.89 \\
Granite 3.0 & 2 & 1.00 & 6.07--9.97 & 6.99 & 6.68 & 0.85 \\
Qwen2.5 & 7 & 1.00 & 2.86--23.75 & 4.11 & 3.53 & 0.72 \\
StableLM-Alpha & 2 & 1.00 & 2.65--3.79 & 2.99 & 2.91 & 0.66 \\
Apertus & 2 & 1.00 & 2.67--6.72 & 3.09 & 2.85 & 0.65 \\
Olmo 3 & 2 & 1.08 & 2.44--4.73 & 2.89 & 2.70 & 0.63 \\
Yi & 2 & 1.00 & 1.81--3.79 & 2.11 & 1.97 & 0.49 \\
Llama 3 herd & 3 & 1.04 & 1.29--6.73 & 1.58 & 1.42 & 0.29 \\
MPT & 2 & 1.00 & 1.20--2.22 & 1.39 & 1.31 & 0.24 \\
Llama 2 & 3 & 1.00 & 1.13--3.04 & 1.45 & 1.29 & 0.23 \\
DeepSeek LLM & 2 & 1.00 & 1.15--2.99 & 1.32 & 1.22 & 0.18 \\
\bottomrule
\end{tabular*}
\par\smallskip
\begin{tablenotes}
The eleven families of the clean inference-demand sample whose members were trained on token budgets within 10 percent
of each other (32 models). Wedges under the reference technology. With a common $D$ and sizes chosen at the margin, the
family's first-order condition in $D$ identifies one number, a weighted average of the members' values of compactness:
$\sum_i\pi_im_{N,i}=\bar w_H-1$ without data costs, where $\pi_i\propto\omega_i/w_i$, $\omega_i$ is member $i$'s share
of family training compute and $\bar w_H=[\sum_i\omega_i/w_i]^{-1}$ is the compute-weighted harmonic mean of the member
wedges, computed from the technology (Proposition~\ref{prop:wedge}(iii); Online Appendix
Proposition~\ref{prop:A-family}); $1-1/\bar w_H=\sum_i\omega_is_i$ is the family share. Member wedges then differ only through the technology, $\ln w_i-\ln w_j=-\alpha(\ln N_i-\ln N_j)$, and are not
revealed-preference objects. If sizes come from a menu, the condition in $D$ says nothing about the members' values of
compactness. In three family problems solved by brute force, the condition holds to within $2.2\times10^{-8}$.
\end{tablenotes}
\begin{tablenotes}[Source]
Model cards and technical reports; authors' calculations.
\end{tablenotes}
\end{table}


% -----------------------------------------------------------------------------------------------------
\subsection{The Cost Side}\label{app:wedge-cost}

FLOPs are not costs. Table~\ref{tab:app-costsens} lets the price of a training FLOP vary with model size (elasticity
$\delta$, for example because utilization changes with size) and serving cost per token scale as $N^\phi$, with a price
$p$ of a serving FLOP relative to a training FLOP (Proposition~\ref{prop:A-wedge}(iii)(f)). The first-order condition is
$w=1+\delta+\phi X_S/X_T$, so the share $s=(w-1-\delta)/(w-1-\delta+\phi)$ does not depend on $p$. It moves little: the
median is 0.72 to 0.77 for $\delta$ between $-0.1$ and $0.1$ with $\phi=1+\delta$, and 0.70 to 0.86 for $\phi$ between 0.5
and 1.25 with $\delta=0$. Token counts do depend on $p$: at the median, planned serving tokens are 9.0 times training
tokens at $p=1$ and 0.9 times at $p=10$, and for Llama~3 8B 17.2 and 1.7 times.

% Online Appendix table relabelled from output/tables/ra2_wedge_costsens.tex (writer appendix_D_F, 2026-09-24).
% [round 3, WP4b, fix list W24(h)] price ratio p -> \varrho (notation decision, fix list section 0.8). The token audit
% leaves every entry unchanged: each column is monotone in w and the median model (w = 3.99) is the same.
\begin{table}[tp]
\centering
\caption{Sensitivity of the Revealed Share to the Cost Side}
\label{tab:app-costsens}
\footnotesize
\setlength{\tabcolsep}{4pt}
\ifdefined\ratabbox\else\newsavebox{\ratabbox}\fi
\sbox{\ratabbox}{%
\begin{tabular}{@{}lrrrrrrrr@{}}
\toprule
 & & & \multicolumn{2}{c}{Share $s$} & \multicolumn{2}{c}{Median $T/D$} & \multicolumn{2}{c}{Llama 3 8B $T/D$} \\
\cmidrule(lr){4-5}\cmidrule(lr){6-7}\cmidrule(l){8-9}
Case & $\delta$ & $\phi$ & median & Llama 3 8B & $\varrho=1$ & $\varrho=10$ & $\varrho=1$ & $\varrho=10$ \\
\midrule
A & $-$0.1 & 0.90 & 0.77 & 0.87 & 10.3 & 1.0 & 19.4 & 1.9 \\
A & 0.0 & 1.00 & 0.75 & 0.85 & 9.0 & 0.9 & 17.2 & 1.7 \\
A & 0.1 & 1.10 & 0.72 & 0.84 & 7.9 & 0.8 & 15.3 & 1.5 \\
B & 0.0 & 0.50 & 0.86 & 0.92 & 17.9 & 1.8 & 34.4 & 3.4 \\
B & 0.0 & 0.75 & 0.80 & 0.88 & 11.9 & 1.2 & 22.9 & 2.3 \\
B & 0.0 & 1.00 & 0.75 & 0.85 & 9.0 & 0.9 & 17.2 & 1.7 \\
B & 0.0 & 1.25 & 0.70 & 0.82 & 7.2 & 0.7 & 13.7 & 1.4 \\
\bottomrule
\end{tabular}%
}
\ifdim\wd\ratabbox>\textwidth\resizebox{\textwidth}{!}{\usebox{\ratabbox}}\else\usebox{\ratabbox}\fi\par
\begin{tablenotes}Lifetime cost $6N^{1+\delta}D+2\varrho N^{\phi}T$ in training-FLOP units at the model's own size: $\delta$ is the elasticity of the training cost of a FLOP with respect to $N$ (model FLOPs utilization falling or rising with size), $\phi$ the elasticity of serving cost per token with respect to $N$, and $\varrho$ the price of a serving FLOP relative to a training FLOP. The first-order condition is $w=1+\delta+\phi X_S/X_T$, with $X_S$ and $X_T$ serving and training expenditure, so $X_S/X_T=(w-1-\delta)/\phi$, $s=(w-1-\delta)/(w-1-\delta+\phi)$ and $T/D=3(w-1-\delta)/(\varrho\phi)$. Case A: $\phi=1+\delta$, so that $w=(1+\delta)(1+\varrho T/3D)$. Case B: $\delta=0$. The share $s$ does not depend on $\varrho$. Clean sample, reference technology; $T/D$ is the serving reading of the wedge (Proposition~\ref{prop:A-wedge}(iii)(a)), which requires that the developer bears the serving cost.\end{tablenotes}
\begin{tablenotes}[Source]Author's calculations.\end{tablenotes}
\end{table}


% -----------------------------------------------------------------------------------------------------
\subsection{Validation and Trends}\label{app:wedge-trends}

\emph{Token volumes.}---OpenRouter's documented interface reports no usage for any model, and its terms of service
(updated 31 August 2026) prohibit scripts that scrape its activity pages, so per-model token volumes cannot be collected.
The developer-level totals of \citet{aubakirova2026state} for November 2024 to November 2025 give an order of magnitude.
One year of OpenRouter traffic equals 4.6 percent of the lifetime tokens that DeepSeek's verified releases imply at $p=1$,
0.7 percent for Moonshot, 0.5 percent for Zhipu, 0.2 percent for Meta, 0.08 percent for Google and 0.04 percent for
Alibaba. The platform carries a small part of a model's lifetime use, and a test of levels is not possible with public
data.

\emph{Derivative repositories.}---From the Hugging Face model tree we count the quantizations, fine-tunes, adapters and
merges derived from each model's base repository, its official post-trained releases and same-run re-releases (74 clean
models from 16 developers; MosaicML's and Cerebras's repositories require authentication). At given compute and release
year, models with higher $M$ have more derivatives: elasticities of 0.61 for quantizations, 0.87 for fine-tunes, 0.90 for
adapters and 0.80 for all derivatives (Table~\ref{tab:app-conduct}, panel~B). The elasticity with respect to the planned
$T/D$ is 0.60 ($p=0.086$), below one. Smaller models of equal compute are adapted more, as planned demand predicts but also
as accessibility predicts, so the counts do not validate levels.

\emph{Trends.}---Figure~\ref{fig:app-trends} plots the median share by release year in the universe of the open--closed
test. Among open-weight models it rises from 0.51 in 2023 to 0.81 in 2024 and 0.85 in 2025. Closed models with disclosed
allocations had negative median shares in 2021--2023 and 0.64 in 2024, from six models; there are too few after 2024 to
plot. The compute-weighted aggregates are in Table~\ref{tab:econ}, panel~C. On the clean sample the aggregate share
rises from 2023 to 2025 under every one of the 32 technologies. Under the developers' own technologies where available it
is 0.27, 0.56 and 0.82 in 2023, 2024 and 2025, against 0.28, 0.58 and 0.83 under the reference. Llama~3.1 405B accounts
for 55 percent of the clean sample's 2024 compute; without it the 2024 share is 0.73.

\begin{figure}[tp]
\centering
\includegraphics[width=0.72\textwidth]{ra2_wedge_trends}
\caption{Revealed Expenditure Shares by Release Year: Open and Closed Models}\label{fig:app-trends}
\begin{figurenotes}
Median of $s=(w-1)/w$ under the reference technology by release year, for open-weight and closed models of the
production-scale universe ($6ND\ge10^{21}$ FLOP) with confident parameter and token counts, excluding mixture-of-experts
and non-transformer models and applying the clean-sample exclusions to verified models (186 models). Shaded: interquartile
range of open-weight models. Numbers: models per cell; cells with fewer than five models are omitted; medians below $-1$
are drawn at $-1$ (open triangles). For developers that serve their models, $s$ is the planned serving share of lifetime
cost; for others, the internalized value of compactness.
\end{figurenotes}
\begin{figurenotes}[Source]
Epoch AI \citep{epochai2026data}; model cards; authors' calculations.
\end{figurenotes}
\end{figure}

\emph{What is not done.}---Four checks remain open. A test of the production approach on a second margin, in the spirit
of \citet{raval2023testing}, such as the sparsity of mixture-of-experts models, the number of key-value heads or the
vocabulary, which trade quality against serving cost at different training costs, has not been run. Token counts are not
normalized by bytes. Durability, the obsolescence of models that shortens the horizon over which serving demand is
planned, is part of $T$ but is not modeled. And token counts reported as corpus sizes, such as the 18 trillion tokens
stated for every Qwen2.5 size, have not been audited against the tokens each model processed.

\clearpage
